Este anejo complementa la descripción metodológica con una lectura de
ingeniería del código. No pretende duplicar el repositorio, sino aislar los
puntos que sostienen la reproducibilidad del trabajo: el contrato de simulación,
la separación entre dinámica, control, telemetría y visualización, y la forma en
que se integra el controlador neuronal dentro del lazo clásico.

\subsubsection*{Organización funcional del paquete}

El paquete \texttt{simulador\_quad} está organizado de forma modular según sus
responsabilidades bajo el directorio \path{src/simulador_quad/}:
\begin{itemize}
    \item \path{dynamics}: Dinámica de cuerpo rígido y actuadores.
    \item \path{core}: Sistemas de referencia, contratos de datos y utilidades comunes.
    \item \path{control}: Implementación de los diferentes controladores.
    \item \path{trajectories}: Generación y gestión de trayectorias de referencia.
    \item \path{telemetry}: Registro estructurado de la telemetría de simulación.
    \item \path{metrics}: Cálculo de métricas de rendimiento y error.
    \item \path{visualization}: Renderizado y salidas gráficas.
\end{itemize}
Los scripts de \path{tools/} orquestan la generación de datasets, bancos de
expertos, entrenamiento neuronal, ejecución masiva y consolidación de
resultados.

Esta separación es relevante porque el trabajo compara controladores sin
cambiar el entorno físico. La campaña puede sustituir el lazo externo por una
red, congelar expertos PD o ejecutar una batería OOD, pero el integrador, el
mixer, los actuadores, las perturbaciones y las métricas siguen siendo comunes.

\subsubsection*{Bucle de simulación multirrate}

El fragmento siguiente resume el núcleo del contrato temporal. En cada ciclo se
comprueban límites de seguridad, se calcula la referencia, se actualiza el
control con retención de orden cero, se aplica el sistema de actuadores, se
registra telemetría a una frecuencia menor y se integra la física mediante
RK4. Esta estructura es la que permite fijar \SI{100}{Hz} para la física,
\SI{50}{Hz} para control y \SI{10}{Hz} para telemetría sin mezclar
responsabilidades.

\begin{lstlisting}[language=Python,caption={Esquema simplificado del bucle de simulación.},label={lst:bucle-simulacion}]
while True:
    is_saturated = any(current_applied.saturation_flags)
    term, reason = self._check_safety_termination(state, is_saturated)
    if term:
        break

    current_ref = trajectory.get_reference_for_state(time_s, state)

    if time_s - last_control_time >= self.control_dt_s:
        obs_pos, obs_vel = self.noise.apply_noise(
            state.position_W_m,
            state.velocity_W_m_s,
        )
        current_control = controller.compute_control(time_s, current_obs, current_ref)
        current_rotor_cmd = self.mixer.compute_rotor_commands(
            current_control.collective_thrust_N,
            current_control.body_moments_Nm,
        )

    current_applied, torque_B, thrust_B = \
        self.actuators.compute_applied_forces(current_rotor_cmd.target_omega_rad_s)

    if time_s - last_telemetry_time >= self.telemetry_dt_s:
        telemetry.append(sample_from_state_control_and_reference(...))

    state = rk4_step(..., thrust_B, torque_B, wind_W_m_s, drag_coefficients)
\end{lstlisting}

\subsubsection*{Integración del controlador neuronal}

El controlador neuronal empleado en la comparación principal no sustituye todo
el piloto automático. La red estima la fuerza deseada del lazo externo en el
sistema mundo y el controlador clásico conserva la conversión fuerza-actitud,
el lazo interno, las saturaciones y el mixer. La inferencia recurrente se hace
con una ventana de longitud fija; por tanto, la red no se ejecuta en modo
\textit{stateful} persistente, sino reconstruyendo la entrada secuencial en cada
paso a partir de la cola FIFO de muestras normalizadas.

\begin{lstlisting}[language=Python,caption={Preparación de entrada para MLP y redes recurrentes.},label={lst:ventana-inferencia}]
x = build_outer_force_min_features_from_observation(obs_state, reference)
x_norm = self.normalizer.normalize_x(torch.tensor(x, dtype=torch.float32))

if self.architecture == "mlp":
    model_input = x_norm.unsqueeze(0)          # [1, input_dim]
else:
    self.window.append(x_norm)
    while len(self.window) < self.sequence_length:
        self.window.appendleft(x_norm)
    model_input = torch.stack(list(self.window)).unsqueeze(0)  # [1, 20, input_dim]

with torch.no_grad():
    y_norm = self.model(model_input).squeeze(0)
\end{lstlisting}

La decisión anterior facilita una interfaz común entre MLP, GRU y LSTM y evita
que el estado oculto dependa de ejecuciones previas. Como contrapartida,
introduce recálculo de las muestras históricas en cada inferencia recurrente.
Por ese motivo, una línea futura razonable es estudiar variantes que conserven
estado interno entre pasos o que actualicen la ventana de forma incremental,
comprobando explícitamente si el cambio altera el comportamiento cerrado.

\subsubsection*{Trazabilidad de salidas}

Cada ejecución escribe \texttt{telemetry.json} y \texttt{metrics.json}. La
telemetría conserva estado, observación, referencia, orden de control, comando
de rotores, actuación aplicada, viento y, cuando procede, fuerzas neuronales
antes y después del recorte. Las métricas agregan RMSE, error máximo, duración,
causa de terminación, saturación, degradación del mixer y metadatos de
procedencia. Esta separación permite reconstruir figuras diagnósticas sin
relanzar la simulación y, a la vez, consolidar campañas completas en CSV.
